\chapter{FEJ-IEKF para Estimación Consistente en Odometría LiDAR-Inercial}
\label{ch:fej-iekf}

\section{Introducción}
\label{sec:fej-introduccion}

El \gls{iekf} es el algoritmo dominante en los sistemas modernos de odometría
\gls{lidar}-inercial, incluyendo FAST-LIO2~\cite{xu2022fastlio2}, sobre el cual se
basa el presente trabajo. Al iterar el punto de linealización hacia la estimación
de máxima a posteriori dentro de cada actualización de medición, el \gls{iekf}
alcanza mayor precisión que el \gls{ekf} estándar con un coste computacional
moderado.

Sin embargo, esta práctica de relinealizar los Jacobianos en el último estado
estimado introduce un artefacto sutil pero significativo: la covarianza estimada
puede reducirse artificialmente en direcciones del espacio de estado que son
fundamentalmente no observables. Este fenómeno fue formalmente analizado por
Huang \emph{et al.}~\cite{huang2010observability} para el caso de \gls{slam}
basado en \gls{ekf}, donde demostraron que la dimensión del subespacio observable
del sistema linealizado puede ser mayor que la del sistema no lineal subyacente.
Como consecuencia, el filtro adquiere información espuria a lo largo de
direcciones donde físicamente no debería existir información alguna.

En el contexto de la odometría \gls{lidar}-inercial, este artefacto se manifiesta
como una covarianza de yaw absoluto que decrece monótonamente a lo largo de la
trayectoria, a pesar de que el yaw absoluto es estructuralmente indeterminable
en ausencia de referencias externas (como \gls{gnss} o magnetómetro). El filtro
se vuelve sobreconfiado, lo que dificulta la integración con sistemas de cierre
de bucle o fusión multi-sensor.

Este capítulo presenta una modificación del \gls{iekf} basada en el enfoque de
Jacobianos de Primeras Estimaciones (\gls{fej}), originalmente propuesto por
Huang \emph{et al.}~\cite{huang2010observability} para el contexto de \gls{slam}
y aquí adaptada al filtro iterativo punto-a-plano de FAST-LIO2. La modificación
consiste en congelar el Jacobiano de medición en el estado propagado al inicio
de cada ciclo de actualización, manteniendo la actualización del residual en cada
iteración para preservar las propiedades de convergencia del \gls{iekf}.

Las contribuciones de este capítulo son:

\begin{itemize}
    \item Una derivación detallada del modelo de estado, dinámica de error y
    Jacobiano de medición para la formulación de FAST-LIO2, verificada contra
    la implementación de referencia.
    \item Un análisis de los errores de linealización del \gls{iekf} estándar,
    distinguiendo entre la variación intra-iteración y la variación entre pasos
    temporales.
    \item La adaptación del \gls{fej} al \gls{iekf} de FAST-LIO2, incluyendo
    consideraciones prácticas como el bloqueo de correspondencias punto-a-plano.
    \item Una validación empírica mediante análisis \gls{nees} contra datos de
    captura de movimiento.
\end{itemize}

El alcance teórico de este capítulo está deliberadamente acotado a lo que puede
demostrarse rigurosamente sin recurrir al análisis del estado conjunto
(drone + mapa) o a herramientas más avanzadas como filtros invariantes en
grupos de Lie. Las extensiones formales se discuten en la
Sección~\ref{sec:fej-trabajo-futuro} como trabajo futuro.

\section{Modelo de Estado y Dinámica}
\label{sec:fej-modelo-estado}

\subsection{Definición del Estado}
\label{sec:fej-estado}

El estado del sistema vive sobre la variedad
$\mathcal{M} = \mathrm{SO}(3) \times \mathbb{R}^{12} \times S^2$:

\begin{equation}
\mathbf{x} = (\mathbf{R}, \mathbf{p}, \mathbf{v}, \mathbf{b}_g, \mathbf{b}_a, \mathbf{g})
\label{eq:state-vector}
\end{equation}

donde:

\begin{itemize}
    \item $\mathbf{R} \in \mathrm{SO}(3)$ es la rotación del marco del cuerpo al
    marco del mundo.
    \item $\mathbf{p} \in \mathbb{R}^3$ es la posición en el marco del mundo.
    \item $\mathbf{v} \in \mathbb{R}^3$ es la velocidad en el marco del mundo.
    \item $\mathbf{b}_g \in \mathbb{R}^3$ es el sesgo del giroscopio.
    \item $\mathbf{b}_a \in \mathbb{R}^3$ es el sesgo del acelerómetro.
    \item $\mathbf{g} \in S^2$ es el vector de gravedad, restringido a la esfera
    de radio $9.81\,\mathrm{m/s^2}$.
\end{itemize}

La dimensión total del espacio tangente del estado es $17$ ($3+3+3+3+3+2$). La
representación de la gravedad en $S^2$ refleja que solamente su dirección es
desconocida, mientras que su magnitud está determinada físicamente.

\begin{remark}
La implementación de FAST-LIO2 incluye adicionalmente los parámetros extrínsecos
LiDAR-IMU $(\mathbf{R}_{LI}, \mathbf{p}_{LI})$ en el estado, los cuales pueden
ser estimados en línea cuando la calibración inicial es incierta. Por claridad
de la exposición, los omitimos del análisis FEJ, ya que son constantes en la
dinámica libre de ruido y no afectan las propiedades de observabilidad
relevantes.
\end{remark}

\subsection{Dinámica en Tiempo Continuo}
\label{sec:fej-dinamica-continua}

La unidad inercial proporciona mediciones de velocidad angular
$\boldsymbol{\omega}_m$ y aceleración lineal $\mathbf{a}_m$, ambas corruptas por
sesgos y ruido blanco gaussiano:

\begin{align}
\boldsymbol{\omega}_m &= \boldsymbol{\omega} + \mathbf{b}_g + \mathbf{n}_g
\label{eq:imu-gyro} \\
\mathbf{a}_m &= \mathbf{R}^\top (\mathbf{a} - \mathbf{g}) + \mathbf{b}_a + \mathbf{n}_a
\label{eq:imu-accel}
\end{align}

donde $\boldsymbol{\omega}$ es la velocidad angular real en el marco del cuerpo,
$\mathbf{a}$ es la aceleración lineal real en el marco del mundo, y
$\mathbf{n}_g, \mathbf{n}_a$ son ruidos blancos gaussianos de media cero.

Despejando las cantidades reales y aplicando las leyes cinemáticas estándar, la
dinámica del estado en tiempo continuo es:

\begin{align}
\dot{\mathbf{R}} &= \mathbf{R} [\boldsymbol{\omega}_m - \mathbf{b}_g - \mathbf{n}_g]_\times
\label{eq:rot-dot} \\
\dot{\mathbf{p}} &= \mathbf{v}
\label{eq:pos-dot} \\
\dot{\mathbf{v}} &= \mathbf{R} (\mathbf{a}_m - \mathbf{b}_a - \mathbf{n}_a) + \mathbf{g}
\label{eq:vel-dot} \\
\dot{\mathbf{b}}_g &= \mathbf{n}_{bg}
\label{eq:bg-dot} \\
\dot{\mathbf{b}}_a &= \mathbf{n}_{ba}
\label{eq:ba-dot} \\
\dot{\mathbf{g}} &= \mathbf{0}
\label{eq:g-dot}
\end{align}

donde $[\cdot]_\times$ denota la matriz antisimétrica del producto vectorial, y
$\mathbf{n}_{bg}, \mathbf{n}_{ba}$ son los ruidos de los sesgos modelados como
caminatas aleatorias.

La ecuación~\eqref{eq:rot-dot} describe la evolución de la rotación bajo
velocidad angular en el marco del cuerpo, con la multiplicación a la derecha
reflejando la convención de perturbación derecha. La
ecuación~\eqref{eq:vel-dot} expresa la segunda ley de Newton en el marco del
mundo: la aceleración del cuerpo es la lectura del acelerómetro, corregida por
sesgo, rotada al marco del mundo, más la gravedad.

\subsection{Estado de Error y Linealización}
\label{sec:fej-estado-error}

El estado vive sobre una variedad, por lo que no podemos restar dos estados
directamente. En su lugar, definimos un \emph{estado de error} en el espacio
tangente.

Para la rotación, usamos la perturbación derecha:

\begin{equation}
\mathbf{R}_{\text{true}} = \mathbf{R}_{\text{est}} \cdot \mathrm{Exp}(\delta\boldsymbol{\theta})
\label{eq:rot-error}
\end{equation}

donde $\delta\boldsymbol{\theta} \in \mathbb{R}^3$ es un vector de rotación pequeño
expresado en el marco del cuerpo, y $\mathrm{Exp}(\cdot)$ es el mapa exponencial
de $\mathfrak{so}(3)$ a $\mathrm{SO}(3)$. Para vectores pequeños:

\begin{equation}
\mathrm{Exp}(\delta\boldsymbol{\theta}) \approx \mathbf{I} + [\delta\boldsymbol{\theta}]_\times
\label{eq:exp-approx}
\end{equation}

Para los componentes del estado en espacios euclidianos, el error es la diferencia
estándar:

\begin{align}
\delta\mathbf{p} &= \mathbf{p}_{\text{true}} - \mathbf{p}_{\text{est}} \\
\delta\mathbf{v} &= \mathbf{v}_{\text{true}} - \mathbf{v}_{\text{est}} \\
\delta\mathbf{b}_g &= \mathbf{b}_{g,\text{true}} - \mathbf{b}_{g,\text{est}} \\
\delta\mathbf{b}_a &= \mathbf{b}_{a,\text{true}} - \mathbf{b}_{a,\text{est}}
\end{align}

Para la gravedad sobre $S^2$, $\delta\mathbf{g} \in \mathbb{R}^2$ vive en el plano
tangente de $S^2$ en el estimado actual.

El estado de error completo es:

\begin{equation}
\delta\mathbf{x} = (\delta\boldsymbol{\theta}, \delta\mathbf{p}, \delta\mathbf{v}, \delta\mathbf{b}_g, \delta\mathbf{b}_a, \delta\mathbf{g}) \in \mathbb{R}^{17}
\label{eq:error-state}
\end{equation}

\subsection{Derivación del Jacobiano F}
\label{sec:fej-jacobiano-f}

La dinámica del estado de error se obtiene linealizando las ecuaciones
\eqref{eq:rot-dot}--\eqref{eq:g-dot} alrededor del estimado actual.

\subsubsection{Derivada del error de rotación}

Partiendo de la definición $\mathbf{R}_{\text{true}} = \mathbf{R}_{\text{est}} \mathrm{Exp}(\delta\boldsymbol{\theta})$
y derivando ambos lados con respecto al tiempo, después de manipulaciones
algebraicas que involucran las identidades de matrices antisimétricas, se obtiene:

\begin{equation}
\dot{\delta\boldsymbol{\theta}} = -[\boldsymbol{\omega}_m - \mathbf{b}_g]_\times \delta\boldsymbol{\theta} - \delta\mathbf{b}_g
\label{eq:dtheta-dot}
\end{equation}

\subsubsection{Derivada del error de posición}

Trivialmente, derivando $\delta\mathbf{p} = \mathbf{p}_{\text{true}} - \mathbf{p}_{\text{est}}$:

\begin{equation}
\dot{\delta\mathbf{p}} = \delta\mathbf{v}
\label{eq:dp-dot}
\end{equation}

\subsubsection{Derivada del error de velocidad}

Tomando la diferencia de la ecuación~\eqref{eq:vel-dot} entre el sistema real y
el estimado, sustituyendo las definiciones del estado de error y aplicando la
identidad del producto vectorial $[\mathbf{a}]_\times \mathbf{b} = -[\mathbf{b}]_\times \mathbf{a}$,
se obtiene:

\begin{equation}
\dot{\delta\mathbf{v}} = -\mathbf{R} [\mathbf{a}_m - \mathbf{b}_a]_\times \delta\boldsymbol{\theta} - \mathbf{R} \delta\mathbf{b}_a + \delta\mathbf{g}
\label{eq:dv-dot}
\end{equation}

Cada término tiene una interpretación física clara: el primero captura el efecto
de un error de rotación que rota incorrectamente el vector de aceleración; el
segundo refleja cómo un error de sesgo del acelerómetro se traduce en un error
de aceleración en el marco del mundo; y el tercero representa la propagación
directa del error de gravedad a la velocidad.

\subsubsection{Derivadas de los sesgos y la gravedad}

En la dinámica libre de ruido, los sesgos y la gravedad son constantes:

\begin{equation}
\dot{\delta\mathbf{b}}_g = \mathbf{0}, \quad \dot{\delta\mathbf{b}}_a = \mathbf{0}, \quad \dot{\delta\mathbf{g}} = \mathbf{0}
\label{eq:bias-grav-dot}
\end{equation}

El ruido aleatorio que afecta a estos componentes entra a través de la matriz de
ruido del proceso $\mathbf{G}$, no a través de $\mathbf{F}$.

\subsubsection{Ensamblando F}

Combinando las ecuaciones \eqref{eq:dtheta-dot}--\eqref{eq:bias-grav-dot}, la
matriz Jacobiana del error es:

\begin{equation}
\mathbf{F} =
\begin{bmatrix}
-[\boldsymbol{\omega}_m - \mathbf{b}_g]_\times & \mathbf{0} & \mathbf{0} & -\mathbf{I} & \mathbf{0} & \mathbf{0} \\
\mathbf{0} & \mathbf{0} & \mathbf{I} & \mathbf{0} & \mathbf{0} & \mathbf{0} \\
-\mathbf{R}[\mathbf{a}_m - \mathbf{b}_a]_\times & \mathbf{0} & \mathbf{0} & \mathbf{0} & -\mathbf{R} & \mathbf{I}_g \\
\mathbf{0} & \mathbf{0} & \mathbf{0} & \mathbf{0} & \mathbf{0} & \mathbf{0} \\
\mathbf{0} & \mathbf{0} & \mathbf{0} & \mathbf{0} & \mathbf{0} & \mathbf{0} \\
\mathbf{0} & \mathbf{0} & \mathbf{0} & \mathbf{0} & \mathbf{0} & \mathbf{0}
\end{bmatrix}
\label{eq:F-matrix}
\end{equation}

donde $\mathbf{I}_g$ es la matriz $3 \times 2$ que mapea perturbaciones del
espacio tangente de $S^2$ al espacio $\mathbb{R}^3$ ambiente.

\subsection{Discretización}
\label{sec:fej-discretizacion}

Para aplicar el filtro a tasas discretas de muestreo del IMU $\Delta t$,
discretizamos $\mathbf{F}$ mediante una expansión de Taylor de primer orden:

\begin{equation}
\boldsymbol{\Phi}_k = \mathbf{I} + \mathbf{F}_k \Delta t + \mathcal{O}(\Delta t^2)
\label{eq:phi-discrete}
\end{equation}

La covarianza del ruido del proceso discreto es:

\begin{equation}
\mathbf{Q}_k = \mathbf{G}_k \mathbf{S}_n \mathbf{G}_k^\top \Delta t
\label{eq:Q-discrete}
\end{equation}

donde $\mathbf{S}_n$ es la matriz diagonal de densidades espectrales de los
ruidos del IMU y de las caminatas aleatorias de los sesgos.

\section{Modelo de Medición}
\label{sec:fej-modelo-medicion}

\subsection{Ecuación de Medición Punto-a-Plano}
\label{sec:fej-punto-plano}

FAST-LIO2 utiliza un modelo de medición \emph{punto-a-plano}: para cada punto
$\mathbf{p}_L$ del escaneo \gls{lidar} expresado en el marco del cuerpo, se
busca su correspondencia con un plano local del mapa, definido por una normal
$\mathbf{n}$ y un punto $\mathbf{q}$ sobre el plano (ambos en el marco del
mundo).

La ecuación implícita de medición es:

\begin{equation}
h(\mathbf{x}) = \mathbf{n}^\top (\mathbf{R} \mathbf{p}_L + \mathbf{p} - \mathbf{q}) = 0
\label{eq:measurement-eq}
\end{equation}

es decir, la distancia con signo del punto transformado al plano es cero. El
filtro busca minimizar este residual sobre todas las correspondencias del
escaneo.

\subsection{Linealización del Modelo de Medición}
\label{sec:fej-linealizacion-medicion}

Perturbamos el estado por $\delta\mathbf{x}$. Solo $\mathbf{R}$ y $\mathbf{p}$
aparecen en $h(\mathbf{x})$, por lo que solamente $\delta\boldsymbol{\theta}$ y
$\delta\mathbf{p}$ producen contribuciones no nulas.

Sustituyendo $\mathbf{R} \to \mathbf{R}(\mathbf{I} + [\delta\boldsymbol{\theta}]_\times)$
y $\mathbf{p} \to \mathbf{p} + \delta\mathbf{p}$:

\begin{align}
h(\mathbf{x} + \delta\mathbf{x}) &\approx \mathbf{n}^\top (\mathbf{R}(\mathbf{I} + [\delta\boldsymbol{\theta}]_\times) \mathbf{p}_L + \mathbf{p} + \delta\mathbf{p} - \mathbf{q}) \nonumber \\
&= h(\mathbf{x}) + \mathbf{n}^\top \mathbf{R} [\delta\boldsymbol{\theta}]_\times \mathbf{p}_L + \mathbf{n}^\top \delta\mathbf{p}
\end{align}

Aplicando la identidad del producto vectorial
$[\delta\boldsymbol{\theta}]_\times \mathbf{p}_L = -[\mathbf{p}_L]_\times \delta\boldsymbol{\theta}$:

\begin{equation}
\delta h = -\mathbf{n}^\top \mathbf{R} [\mathbf{p}_L]_\times \delta\boldsymbol{\theta} + \mathbf{n}^\top \delta\mathbf{p}
\label{eq:dh-linearized}
\end{equation}

\subsection{Jacobiano de Medición H}
\label{sec:fej-jacobiano-h}

Para una correspondencia $i$, el Jacobiano de medición es un vector fila
$1 \times 17$:

\begin{equation}
\mathbf{H}_i =
\begin{bmatrix}
-\mathbf{n}_i^\top \mathbf{R} [\mathbf{p}_{L,i}]_\times & \mathbf{n}_i^\top & \mathbf{0} & \mathbf{0} & \mathbf{0} & \mathbf{0}
\end{bmatrix}
\label{eq:H-row}
\end{equation}

Para $N$ correspondencias en un escaneo, las filas se apilan formando una matriz
$\mathbf{H}$ de tamaño $N \times 17$.

El ruido de medición se modela como gaussiano de media cero con varianza
$\sigma^2$ por residual:

\begin{equation}
\mathbf{R} = \sigma^2 \mathbf{I}_N
\label{eq:meas-noise}
\end{equation}

\subsection{Observaciones Clave}
\label{sec:fej-observaciones-h}

Tres observaciones de la estructura de $\mathbf{H}$ son fundamentales para el
análisis de observabilidad:

\begin{enumerate}
    \item Las mediciones \gls{lidar} acoplan directamente solo con la rotación
    y la posición. Velocidad, sesgos y gravedad son observables únicamente de
    forma indirecta a través de la dinámica del sistema.
    \item El bloque de posición $\mathbf{n}_i^\top$ es independiente del estado.
    Solo el bloque de rotación $-\mathbf{n}_i^\top \mathbf{R} [\mathbf{p}_{L,i}]_\times$
    cambia cuando se modifica el punto de linealización.
    \item Cuando el \gls{iekf} itera y reevalúa $\mathbf{H}$ en diferentes
    estimados de $\mathbf{R}$, el bloque de rotación rota con el estimado. Esta
    es la fuente de errores de linealización intra-iteración que el \gls{fej}
    elimina.
\end{enumerate}

\section{Observabilidad de la Estimación LiDAR-Inercial}
\label{sec:fej-observabilidad}

\subsection{Sistema Linealizado en Tiempo Discreto}
\label{sec:fej-sistema-discreto}

Combinando la propagación discretizada y el modelo de medición linealizado, el
sistema de error en tiempo discreto es:

\begin{align}
\delta\mathbf{x}_{k+1} &= \boldsymbol{\Phi}_k \delta\mathbf{x}_k + \mathbf{w}_k \\
\mathbf{z}_k &= \mathbf{H}_k \delta\mathbf{x}_k + \mathbf{v}_k
\end{align}

donde $\mathbf{w}_k$ y $\mathbf{v}_k$ son ruidos gaussianos de media cero con
covarianzas $\mathbf{Q}_k$ y $\mathbf{R}$ respectivamente.

\subsection{Matriz de Observabilidad}
\label{sec:fej-matriz-observabilidad}

Para una ventana de $k+1$ pasos de tiempo, la matriz de observabilidad del
sistema linealizado es:

\begin{equation}
\mathbf{O}_k =
\begin{bmatrix}
\mathbf{H}_0 \\
\mathbf{H}_1 \boldsymbol{\Phi}_0 \\
\mathbf{H}_2 \boldsymbol{\Phi}_1 \boldsymbol{\Phi}_0 \\
\vdots \\
\mathbf{H}_k \boldsymbol{\Phi}_{k-1} \cdots \boldsymbol{\Phi}_0
\end{bmatrix}
\label{eq:obs-matrix}
\end{equation}

Una dirección $\mathbf{n}$ del espacio de estado es no observable si y solo si
$\mathbf{O}_k \mathbf{n} = \mathbf{0}$ para cualquier ventana de observación.
El espacio nulo de $\mathbf{O}_k$ es el \emph{subespacio no observable} del
sistema linealizado.

\subsection{Simetrías Físicas del Sistema}
\label{sec:fej-simetrias}

El sistema \gls{lidar}-inercial posee una única simetría física que produce
direcciones no observables: la \emph{rotación global de yaw alrededor de la
gravedad}.

Específicamente:

\begin{itemize}
    \item \textbf{Roll y pitch son observables}: el acelerómetro detecta el
    vector de gravedad, que actúa como una referencia vertical absoluta.
    \item \textbf{Posición es observable}: el marco del mundo está anclado por
    el primer escaneo \gls{lidar}, y los escaneos subsiguientes se registran
    contra este mapa anclado.
    \item \textbf{Yaw absoluto no es observable}: rotar el sistema entero
    (drone + mapa) alrededor del eje de gravedad no cambia ninguna medición.
    El acelerómetro no lo detecta porque la gravedad es invariante bajo esta
    rotación. El \gls{lidar} no lo detecta porque el mapa rota con el drone.
    El giroscopio mide velocidad angular, no orientación absoluta, por lo que
    no proporciona referencia.
\end{itemize}

\subsection{Resultado de SLAM y Herencia a LIO}
\label{sec:fej-herencia-slam}

Una prueba formal de la no observabilidad del yaw fue establecida por Huang
\emph{et al.}~\cite{huang2010observability} para el caso de \gls{ekf}-\gls{slam}
en 2D, donde demostraron que la matriz de observabilidad tiene un espacio nulo
de dimensión 3 (2 dimensiones de traslación global + 1 dimensión de rotación
global). La extensión a 3D produce un espacio nulo de dimensión 4 (3D
traslación + 1D yaw).

En nuestra formulación de \gls{lio}, el mapa no forma parte del vector de
estado del filtro. En su lugar, se construye incrementalmente en una estructura
\gls{ikdtree} a partir de los estimados del estado del drone. El mapa es por
tanto una función implícita de la trayectoria estimada.

La no observabilidad del yaw global, establecida formalmente en el contexto de
\gls{slam}, se hereda a la formulación \gls{lio} a través de esta dependencia
implícita: cualquier deriva en el yaw absoluto del drone se refleja en una
rotación correspondiente del mapa construido, dejando inalterados todos los
residuales de medición. Aunque el estado del drone es localmente observable
por escaneo (cada escaneo proporciona restricciones contra el mapa actual), el
\emph{yaw absoluto del sistema conjunto (drone + mapa implícito)} no tiene
referencia externa.

\subsection{Generador del Yaw en el Espacio Tangente}
\label{sec:fej-generador-yaw}

La rotación infinitesimal de yaw alrededor de la gravedad actúa sobre el estado
de la siguiente manera:

\begin{itemize}
    \item Rotación: $\delta\boldsymbol{\theta}_{\text{cuerpo}} = \psi \mathbf{R}^\top \hat{\mathbf{g}}$
    \item Posición: $\delta\mathbf{p} = \psi \hat{\mathbf{g}} \times \mathbf{p}$
    \item Velocidad: $\delta\mathbf{v} = \psi \hat{\mathbf{g}} \times \mathbf{v}$
    \item Sesgos y gravedad: invariantes
\end{itemize}

donde $\hat{\mathbf{g}}$ es el vector unitario de gravedad en el marco del
mundo, y $\psi$ es el ángulo infinitesimal de rotación.

El vector tangente correspondiente para $\psi = 1$ es:

\begin{equation}
\mathbf{n}_{\text{yaw}} =
\begin{bmatrix}
\mathbf{R}^\top \hat{\mathbf{g}} \\
\hat{\mathbf{g}} \times \mathbf{p} \\
\hat{\mathbf{g}} \times \mathbf{v} \\
\mathbf{0} \\
\mathbf{0} \\
\mathbf{0}
\end{bmatrix}
\label{eq:n-yaw}
\end{equation}

\begin{remark}
La verificación directa de $\mathbf{H} \mathbf{n}_{\text{yaw}} = \mathbf{0}$ en
la formulación \gls{lio} (sin el mapa en el estado) produce un residual no nulo
que se cancelaría únicamente si los puntos del mapa también rotaran. La prueba
formal para \gls{lio} requiere o bien el análisis del estado conjunto, o bien
herramientas de invariancia en grupos de Lie. Estas extensiones se discuten en
la Sección~\ref{sec:fej-trabajo-futuro}. Para el alcance de este capítulo,
adoptamos el resultado de \gls{slam} de Huang \emph{et al.} y argumentamos la
herencia a \gls{lio} a través de la dependencia implícita del mapa.
\end{remark}

\section{Inconsistencia del IEKF Estándar}
\label{sec:fej-inconsistencia}

\subsection{Dependencia de H del Punto de Linealización}
\label{sec:fej-dependencia-H}

Recordando la ecuación~\eqref{eq:H-row}, los dos bloques no nulos de $\mathbf{H}_i$
son:

\begin{itemize}
    \item Bloque de posición: $\mathbf{n}_i^\top$, que depende solamente de la
    normal del plano y no del estado.
    \item Bloque de rotación: $-\mathbf{n}_i^\top \mathbf{R} [\mathbf{p}_{L,i}]_\times$,
    que depende de $\mathbf{R}$, el punto de linealización.
\end{itemize}

Cuando el \gls{iekf} itera, el bloque de rotación cambia entre iteraciones porque
$\mathbf{R}$ cambia. Cuando el \gls{ekf} se ejecuta a través de múltiples pasos
de tiempo, $\mathbf{R}$ también cambia entre pasos. Cada matriz $\mathbf{H}$ está
por tanto atada a un punto de linealización específico.

\subsection{Variación Intra-Iteración}
\label{sec:fej-variacion-intra}

Dentro de una sola actualización en el tiempo $k$, el \gls{iekf} realiza
múltiples iteraciones:

\begin{itemize}
    \item Iteración 0: estimado del estado $\mathbf{x}_k^{(0)} = \mathbf{x}_{k|k-1}$ (el a priori propagado)
    \item Iteración 1: estimado del estado $\mathbf{x}_k^{(1)} =$ resultado de un paso de actualización del EKF
    \item Iteración 2: estimado del estado $\mathbf{x}_k^{(2)} =$ resultado de dos pasos
    \item $\dots$
\end{itemize}

En cada iteración, la rotación $\mathbf{R}$ cambia ligeramente (típicamente por
miliradianes). El Jacobiano de medición correspondiente cambia:

\begin{equation}
\mathbf{H}_i^{(j)} - \mathbf{H}_i^{(j-1)} =
\begin{bmatrix}
-\mathbf{n}_i^\top (\mathbf{R}^{(j)} - \mathbf{R}^{(j-1)}) [\mathbf{p}_{L,i}]_\times & \mathbf{0} & \cdots
\end{bmatrix}
\label{eq:H-variation}
\end{equation}

Esta es una perturbación pequeña, pero no nula. El \gls{iekf} estándar utiliza
$\mathbf{H}_i^{(j)}$ en la iteración $j$ para calcular la ganancia de Kalman y
la actualización de covarianza, por lo que cada iteración refleja una
linealización distinta.

\subsection{Variación Entre Pasos Temporales}
\label{sec:fej-variacion-entre-pasos}

Incluso dejando de lado las iteraciones del \gls{iekf}, el punto de
linealización cambia entre pasos consecutivos. El a priori propagado
$\mathbf{R}_{k|k-1}$ difiere del a posteriori previo $\mathbf{R}_{k-1|k-1}$
por el paso de integración del \gls{imu}, y después de la actualización,
$\mathbf{R}_{k|k}$ difiere de $\mathbf{R}_{k|k-1}$ por la corrección del
\gls{ekf}.

Apilando mediciones sobre una ventana de pasos temporales, la matriz de
observabilidad es:

\begin{equation}
\mathbf{O} =
\begin{bmatrix}
\mathbf{H}_0(\mathbf{R}_{0|0}) \\
\mathbf{H}_1(\mathbf{R}_{1|0}) \boldsymbol{\Phi}_0(\mathbf{R}_{0|0}) \\
\mathbf{H}_2(\mathbf{R}_{2|1}) \boldsymbol{\Phi}_1(\mathbf{R}_{1|1}) \boldsymbol{\Phi}_0(\mathbf{R}_{0|0}) \\
\vdots
\end{bmatrix}
\label{eq:O-shifting}
\end{equation}

Cada $\mathbf{H}_k$ y cada $\boldsymbol{\Phi}_k$ se evalúan en un estado
distinto. No hay un punto de linealización consistente a través de la matriz,
por lo que su espacio nulo no es directamente comparable a ningún espacio nulo
``verdadero'' derivado de un único punto de evaluación.

\subsection{Conexión con el Resultado de Consistencia de SLAM}
\label{sec:fej-conexion-slam}

Para el caso de \gls{slam}, Huang \emph{et al.}~\cite{huang2010observability}
demostraron que cuando los Jacobianos del \gls{ekf} se evalúan en los últimos
estimados del estado en cada paso de tiempo, el sistema linealizado de error
tiene un subespacio observable de dimensión mayor que la del sistema no lineal
real. Como consecuencia, los estimados de covarianza se reducen en direcciones
del espacio de estado donde no hay información disponible.

El mecanismo es análogo en nuestro caso: los puntos de linealización cambiantes
inyectan información espuria. En \gls{slam}, esta inconsistencia se manifiesta
como una reducción de 1 dimensión en el espacio nulo (el yaw se vuelve
``espuriamente observable''). En nuestro \gls{lio}, se manifiesta como una
covarianza sobreconfiada en la misma dirección, aunque la estructura del estado
de \gls{lio} difiere.

\subsection{Manifestación Empírica en LIO}
\label{sec:fej-manifestacion-empirica}

En la formulación \gls{lio}, la inconsistencia se manifiesta de la siguiente
manera:

\begin{enumerate}
    \item Cada escaneo \gls{lidar} proporciona restricciones geométricas
    relativas al \emph{mapa actual}, que fue construido a partir de estimados
    previos del estado.
    \item El \gls{ekf} trata estas restricciones como mediciones independientes
    de una referencia absoluta, ganando información de Fisher en cada paso.
    \item La covarianza del yaw se reduce en consecuencia, aunque el yaw
    absoluto es fundamentalmente indeterminable (el mapa deriva con el drone).
    \item A lo largo de una trayectoria larga, la incertidumbre publicada del
    yaw se vuelve mucho menor que el error real contra la verdad de campo.
\end{enumerate}

Esta inconsistencia puede medirse directamente mediante el \gls{nees}. Para un
estimador consistente, el \gls{nees} debe ser aproximadamente igual a la
dimensión del estado, con límites de confianza chi-cuadrado. Para el \gls{iekf}
estándar en \gls{lio}, el \gls{nees} excede el límite superior chi-cuadrado
después de unos pocos segundos de operación, indicando que la covarianza
publicada es demasiado pequeña en relación con el error real.

\subsection{Alcance del Análisis}
\label{sec:fej-alcance}

Para evitar afirmaciones que excedan el rigor formal disponible en este
capítulo, distinguimos claramente lo que podemos y no podemos demostrar.

\textbf{Podemos afirmar:}

\begin{itemize}
    \item El \gls{iekf} estándar relinealiza $\mathbf{H}$ en cada iteración y
    en cada paso temporal.
    \item Esta relinealización introduce errores pequeños pero acumulativos.
    \item Los errores se manifiestan como estimaciones de covarianza
    sobreconfiadas, particularmente en direcciones relacionadas con la simetría
    de yaw global.
    \item Esto es una instancia del mismo mecanismo que Huang \emph{et al.}
    analizaron formalmente para \gls{slam}.
    \item La prueba empírica \gls{nees} detecta esta inconsistencia en \gls{lio}.
\end{itemize}

\textbf{No podemos afirmar} (sin el análisis del estado conjunto):

\begin{itemize}
    \item Una prueba formal de que el \gls{iekf} de \gls{lio} viola una
    condición específica de dimensión del espacio nulo.
    \item Una expresión analítica para la tasa de ganancia de información falsa
    en \gls{lio}.
    \item Una correspondencia directa entre la inconsistencia de \gls{lio} y la
    de \gls{slam} a nivel de dimensiones.
\end{itemize}

Estas afirmaciones más fuertes requieren o bien el análisis del estado conjunto
(drone + mapa) o bien el enfoque de invariancia en grupos de Lie, ambos
diferidos a trabajo futuro.

\section{El Algoritmo FEJ-IEKF}
\label{sec:fej-algoritmo}

\subsection{Modificación FEJ}
\label{sec:fej-modificacion}

La idea central del \gls{fej} es congelar el Jacobiano de medición $\mathbf{H}$
en el estado a priori propagado al inicio de cada ciclo de actualización, en
lugar de reevaluarlo en cada iteración. Esto elimina la fuente intra-iteración
de errores de linealización identificada en la
Sección~\ref{sec:fej-variacion-intra}.

\subsection{Ecuaciones de Actualización}
\label{sec:fej-ecuaciones-actualizacion}

\textbf{Inicialización ($j = 0$):}

\begin{equation}
\mathbf{H}_k^{\text{FEJ}} = \left.\frac{\partial h}{\partial \delta\mathbf{x}}\right|_{\mathbf{x}_{k|k-1}}
\label{eq:H-fej}
\end{equation}

\begin{equation}
\mathbf{K}_k = \mathbf{P}_{k|k-1} (\mathbf{H}_k^{\text{FEJ}})^\top \left(\mathbf{H}_k^{\text{FEJ}} \mathbf{P}_{k|k-1} (\mathbf{H}_k^{\text{FEJ}})^\top + \mathbf{R}\right)^{-1}
\label{eq:K-fej}
\end{equation}

\textbf{Iteración $j$:}

\begin{enumerate}
    \item Calcular el residual en el estimado actual:
    \begin{equation}
    \mathbf{z}_k^{(j)} = h_{\text{actual}} - h(\mathbf{x}_k^{(j)})
    \label{eq:residual-fej}
    \end{equation}

    \item Calcular la actualización del estado usando $\mathbf{H}^{\text{FEJ}}$ y $\mathbf{K}_k$ fijos:
    \begin{equation}
    \delta\mathbf{x}_k^{(j)} = \mathbf{K}_k \left(\mathbf{z}_k^{(j)} - \mathbf{H}_k^{\text{FEJ}} (\mathbf{x}_k^{(j)} \boxminus \mathbf{x}_{k|k-1})\right)
    \label{eq:state-update-fej}
    \end{equation}

    \item Aplicar la actualización sobre la variedad:
    \begin{equation}
    \mathbf{x}_k^{(j+1)} = \mathbf{x}_k^{(j)} \boxplus \delta\mathbf{x}_k^{(j)}
    \label{eq:manifold-update-fej}
    \end{equation}

    \item Verificar convergencia: si $\|\delta\mathbf{x}_k^{(j)}\| < \epsilon$, detener.
\end{enumerate}

\textbf{Actualización de covarianza tras convergencia:}

\begin{equation}
\mathbf{P}_{k|k} = (\mathbf{I} - \mathbf{K}_k \mathbf{H}_k^{\text{FEJ}}) \mathbf{P}_{k|k-1}
\label{eq:cov-update-fej}
\end{equation}

\subsection{Comparación con el IEKF Estándar}
\label{sec:fej-comparacion}

\begin{table}[t]
\centering
\caption{Comparación entre IEKF estándar y FEJ-IEKF}
\label{tab:iekf-fej-comparison}
\begin{tabular}{lll}
\hline
\textbf{Elemento} & \textbf{IEKF Estándar} & \textbf{FEJ-IEKF} \\
\hline
Residual $\mathbf{z}_k^{(j)}$ & Recalculado en cada iter. & Recalculado en cada iter. \\
Jacobiano $\mathbf{H}_k$ & Recalculado en cada iter. & Congelado en iter. 0 \\
Ganancia $\mathbf{K}_k$ & Recalculada en cada iter. & Calculada una vez en iter. 0 \\
Actualización del estado & Usa $\mathbf{H}^{(j)}$, $\mathbf{K}^{(j)}$ & Usa $\mathbf{H}^{\text{FEJ}}$, $\mathbf{K}$ \\
Actualización de covarianza & Usa $\mathbf{H}^{(*)}$ (final) & Usa $\mathbf{H}^{\text{FEJ}}$ \\
\hline
\end{tabular}
\end{table}

El comportamiento de convergencia se preserva porque $\mathbf{z}_k^{(j)}$ se
recalcula (es el residual el que dirige la convergencia), mientras que la
propiedad de consistencia se gana porque $\mathbf{H}_k^{\text{FEJ}}$ está
congelado (es el Jacobiano el que determina el punto de linealización para la
ganancia y la covarianza).

\subsection{Lo Que FEJ Corrige y Lo Que No}
\label{sec:fej-alcance-correccion}

\textbf{FEJ corrige:}

\begin{itemize}
    \item Elimina la deriva intra-iteración en $\mathbf{H}$, $\mathbf{K}$ y la
    actualización de covarianza.
    \item La ganancia de Kalman refleja una linealización única y consistente
    dentro de un ciclo de actualización.
    \item La reducción de covarianza en cada actualización se calcula en un
    único punto de linealización consistente.
\end{itemize}

\textbf{FEJ no corrige:}

\begin{itemize}
    \item La variación entre pasos temporales identificada en la
    Sección~\ref{sec:fej-variacion-entre-pasos} no es abordada por \gls{fej}.
    A través de pasos consecutivos, el punto de linealización todavía cambia
    ($\mathbf{x}_{k|k-1}$ difiere de $\mathbf{x}_{k-1|k-1}$).
    \item La covarianza no crece monótonamente con el tiempo. \gls{fej}
    elimina la reducción artificial de la covarianza, pero no añade
    incertidumbre que refleje la deriva acumulada a lo largo de la
    trayectoria. La covarianza alcanza un mínimo no degenerado y se estabiliza
    en él.
    \item La covarianza acotada resultante refleja correctamente la
    incertidumbre local instantánea, pero no la incertidumbre absoluta del
    drone respecto al sistema de referencia inicial. Para aplicaciones que
    requieren una estimación honesta de la deriva acumulada (cierre de bucle,
    fusión \gls{gnss}), se requieren mecanismos complementarios discutidos en
    la Sección~\ref{sec:fej-trabajo-futuro}.
\end{itemize}

\section{Implementación}
\label{sec:fej-implementacion}

La implementación del \gls{fej}-\gls{iekf} en este proyecto extiende el conjunto
de herramientas IKFoM utilizado por FAST-LIO2. La nueva función
\texttt{update\_iterated\_dyn\_share\_modified\_fej} en
\texttt{esekfom.hpp} sigue exactamente las ecuaciones presentadas en la
Sección~\ref{sec:fej-ecuaciones-actualizacion}.

\subsection{Bloqueo de Correspondencias Punto-a-Plano}
\label{sec:fej-bloqueo-puntos}

Una consideración práctica importante es que la función \texttt{h\_share\_model}
de FAST-LIO2 recalcula las correspondencias punto-a-plano en cada iteración del
\gls{iekf}. Si las correspondencias cambian entre iteraciones, las dimensiones
de $\mathbf{H}_k^{\text{FEJ}}$ pueden no coincidir con las dimensiones del
residual en iteraciones posteriores, invalidando el Jacobiano congelado.

La implementación maneja esto con un flag \texttt{fej\_lock\_points} que indica
a \texttt{h\_share\_model} mantener fijas las correspondencias durante todas las
iteraciones de un ciclo de actualización. Las correspondencias se calculan una
vez en la iteración 0 y se reutilizan en las iteraciones subsiguientes.

\subsection{Manejo de Cambios de Dimensión}
\label{sec:fej-cambio-dimension}

Como mecanismo de respaldo, si las correspondencias cambian inesperadamente
(por ejemplo, debido a casos límite numéricos), la implementación recaptura
$\mathbf{H}_k^{\text{FEJ}}$ en la iteración donde ocurre el cambio, en lugar de
fallar. Esto garantiza la robustez en condiciones edge-case sin sacrificar la
propiedad de consistencia en el caso común.

\subsection{Coste Computacional}
\label{sec:fej-coste-computacional}

Notablemente, el \gls{fej}-\gls{iekf} tiene un coste computacional ligeramente
\emph{menor} que el \gls{iekf} estándar, ya que:

\begin{itemize}
    \item El Jacobiano $\mathbf{H}_k^{\text{FEJ}}$ se calcula una sola vez en
    lugar de en cada iteración.
    \item La ganancia de Kalman $\mathbf{K}_k$ se calcula una sola vez en lugar
    de en cada iteración.
    \item Solamente el residual $\mathbf{z}_k^{(j)}$ se recalcula en cada
    iteración, lo cual es mucho más barato que recalcular $\mathbf{H}$.
\end{itemize}

En la práctica, la diferencia es mínima (unos pocos por ciento) porque el
cálculo del residual domina de todos modos. Pero no hay sobrecarga al usar
\gls{fej}, solo ahorro.

\section{Validación Empírica}
\label{sec:fej-validacion}

La validación experimental del \gls{fej}-\gls{iekf} se presenta en detalle en
el Capítulo~\ref{ch:experimentos}. Aquí resumimos brevemente la metodología y
los resultados clave.

\subsection{Análisis NEES}
\label{sec:fej-nees}

La métrica principal es el \gls{nees}, calculado contra verdad de campo
proporcionada por un sistema de captura de movimiento OptiTrack en un recinto
de $9 \times 9\,\mathrm{m}$. Para un estimador consistente, el \gls{nees} debe
estar dentro de los límites chi-cuadrado correspondientes a la dimensión del
estado.

Los resultados muestran que:

\begin{itemize}
    \item El \gls{iekf} estándar produce \gls{nees} que excede el límite
    superior chi-cuadrado después de unos pocos segundos de operación, lo
    cual indica covarianza sobreconfiada.
    \item El \gls{fej}-\gls{iekf} mantiene el \gls{nees} dentro de los límites
    durante un período significativamente mayor, debido a la eliminación de
    la fuente intra-iteración de información espuria.
    \item En trayectorias suficientemente largas, el \gls{nees} del
    \gls{fej}-\gls{iekf} eventualmente excede los límites a medida que la
    deriva acumulada supera la covarianza acotada del filtro. Abordar esta
    limitación requiere mecanismos adicionales que se discuten en la
    Sección~\ref{sec:fej-trabajo-futuro}.
\end{itemize}

\subsection{Evolución de la Covarianza de Yaw}
\label{sec:fej-yaw-cov}

Una segunda métrica relevante es la evolución de la covarianza del yaw a lo
largo del tiempo. Para el \gls{iekf} estándar, la covarianza del yaw decrece
monótonamente, tendiendo asintóticamente a cero. Esto no es físicamente
justificado, ya que el yaw absoluto es estructuralmente indeterminable en
ausencia de referencias externas, y reflejarlo con una covarianza decreciente
constituye una sobreconfianza.

Para el \gls{fej}-\gls{iekf}, la covarianza del yaw deja de decrecer hacia
cero y se estabiliza en un valor mínimo no degenerado. Esta es la mejora
directa que aporta el congelamiento del Jacobiano: la fuente intra-iteración
de información espuria se elimina, y la covarianza ya no se reduce
artificialmente.

Sin embargo, debe notarse que el \gls{fej}-\gls{iekf} por sí solo no produce
una covarianza monótonamente creciente. Cada nueva medición sigue
proporcionando información real en las direcciones localmente observables,
pero el mecanismo \gls{fej} no añade incertidumbre con el tiempo. El resultado
es una covarianza que permanece acotada en un valor mínimo, en lugar de
colapsar a cero.

Esto representa una mejora significativa de consistencia local respecto al
\gls{iekf} estándar, pero no captura la incertidumbre acumulada de deriva a
lo largo de trayectorias largas. La covarianza publicada no refleja el hecho
de que el error absoluto del drone respecto al sistema de referencia inicial
crece con la distancia recorrida. Esta limitación se discute como trabajo
futuro en la Sección~\ref{sec:fej-trabajo-futuro}.

\section{Limitaciones y Trabajo Futuro}
\label{sec:fej-trabajo-futuro}

\subsection{Limitaciones del Análisis Presentado}

\begin{enumerate}
    \item \textbf{Falta de prueba formal de observabilidad para LIO}: hemos
    citado el resultado de \gls{slam} de Huang \emph{et al.} y argumentado la
    herencia a \gls{lio} mediante la dependencia implícita del mapa, pero no
    hemos derivado formalmente el espacio nulo del sistema \gls{lio} desde
    primeros principios.
    \item \textbf{Variación entre pasos temporales no abordada}: el \gls{fej}
    soluciona la variación intra-iteración pero no la variación entre pasos
    consecutivos. La covarianza producida por \gls{fej}-\gls{iekf} es
    localmente consistente pero no captura la incertidumbre acumulada de
    deriva a lo largo de trayectorias largas.
    \item \textbf{Covarianza no monótonamente creciente}: \gls{fej} elimina
    la reducción artificial de la covarianza, pero no añade incertidumbre con
    el tiempo. Para aplicaciones que requieren una estimación honesta de la
    incertidumbre absoluta acumulada, esto es una limitación que debe
    abordarse con mecanismos complementarios.
    \item \textbf{Validación restringida a entornos controlados}: la validación
    \gls{nees} requiere verdad de campo precisa, disponible solamente en el
    recinto de captura de movimiento. La extrapolación a entornos exteriores
    o de larga distancia se basa en métricas indirectas como el error de cierre
    de bucle.
\end{enumerate}

\subsection{Direcciones de Trabajo Futuro}

\subsubsection{Prueba Formal de Observabilidad para LIO}

Tres enfoques pueden formalizar el subespacio no observable del \gls{lio}
\gls{iekf} sin depender del análogo de \gls{slam}:

\begin{enumerate}
    \item \textbf{Análisis del estado conjunto}: aumentar el estado con puntos
    del mapa (o surfels) de forma explícita, aplicar la derivación del espacio
    nulo de \gls{slam}, y mostrar que el resultado se reduce correctamente
    cuando el mapa se marginaliza.
    \item \textbf{Invariancia en grupos de Lie}: enmarcar la simetría de yaw
    como una acción de grupo sobre el estado, mostrar que la dinámica es
    equivariante y las mediciones son invariantes bajo la acción conjunta
    (estado + mapa), y derivar el generador del espacio tangente sin necesidad
    del mapa en el estado.
    \item \textbf{Argumento de marginalización}: calcular la matriz de
    información conjunta sobre (drone, mapa), marginalizar el mapa, y mostrar
    que la matriz de información del drone resultante tiene un valor propio
    cero en la dirección del yaw.
\end{enumerate}

\subsubsection{OC-IEKF}

El enfoque \gls{ocekf} de Huang \emph{et al.} aborda la variación entre pasos
temporales que el \gls{fej} no corrige, imponiendo restricciones adicionales
sobre $\boldsymbol{\Phi}$ y $\mathbf{H}$ más allá de simplemente congelar el
punto de linealización por actualización. Su adaptación al \gls{iekf} de
\gls{lio} requeriría una derivación desde primeros principios en el marco de
estado conjunto, constituyendo una contribución sustancial adecuada para una
publicación de revista.

\subsubsection{Integración con Filtros Invariantes}

Una comparación con el filtro \gls{ekf} invariante (basado en estructura de
grupos de Lie) revelaría si el \gls{fej} alcanza una consistencia similar a
través de un mecanismo más simple, o si los filtros invariantes ofrecen
ventajas adicionales en términos de robustez o tasa de convergencia.

\subsubsection{Extensiones a Otras Modalidades}

El enfoque \gls{fej} debería generalizarse naturalmente a:

\begin{itemize}
    \item Odometría visual-inercial (mediciones de reproyección de puntos)
    \item Odometría inercial con rango (balizas \gls{uwb})
    \item Fusión multi-modal
\end{itemize}

Cada modalidad tiene su propia estructura de $\mathbf{H}$ y propiedades de
observabilidad, pero el principio general de congelar el punto de linealización
por ciclo de actualización aplica universalmente.

\section{Conclusiones del Capítulo}
\label{sec:fej-conclusiones}

En este capítulo hemos presentado una adaptación del enfoque de Jacobianos de
Primeras Estimaciones al \gls{iekf} de FAST-LIO2 para mejorar la consistencia
de las estimaciones de covarianza en odometría \gls{lidar}-inercial. La
modificación es algorítmicamente simple pero conceptualmente significativa:
congelar el Jacobiano de medición $\mathbf{H}$ en el estado a priori propagado
al inicio de cada ciclo de actualización, mientras se mantiene la actualización
del residual en cada iteración para preservar las propiedades de convergencia
del \gls{iekf}.

El análisis presentado muestra que esta modificación elimina la fuente
intra-iteración de errores de linealización del \gls{iekf} estándar, los cuales
son la manifestación local del mismo mecanismo que Huang \emph{et al.}
identificaron formalmente para el caso de \gls{slam}. La validación empírica
mediante \gls{nees} contra verdad de campo confirma que la covarianza producida
por el \gls{fej}-\gls{iekf} no se reduce artificialmente hacia cero como en el
caso del \gls{iekf} estándar, sino que se estabiliza en un valor mínimo no
degenerado que refleja correctamente la incertidumbre local instantánea.

Las limitaciones formales del análisis (la imposibilidad de derivar
directamente el espacio nulo del \gls{lio} sin recurrir al análisis del estado
conjunto, y la variación entre pasos temporales que el \gls{fej} no aborda) se
identifican explícitamente y se proponen como direcciones de trabajo futuro.
La contribución de este capítulo no es la generalización formal del resultado
de Huang \emph{et al.} a \gls{lio}, sino su \emph{adaptación práctica} al
filtro de FAST-LIO2 con consideraciones de implementación específicas (bloqueo
de correspondencias, manejo de cambios de dimensión) y validación empírica en
una plataforma multirrotor real.
